\section{Physical phase branch: induction, proper time, Unruh, Aharonov--Bohm, Berry, Euler}

\subsection{Faraday and Maxwell as the classical control layer}

For a coil with \(N\) turns and magnetic flux \(\Phi_B\), Faraday's law gives
\begin{equation}
\boxed{\mathcal E=-N\frac{d\Phi_B}{dt}.}
\end{equation}
For the idealised rotating-flux control
\begin{equation}
\Phi_B(t)=BA\cos(\omega t+\phi),
\end{equation}
one obtains
\begin{equation}
\boxed{\mathcal E(t)=NAB\omega\sin(\omega t+\phi).}
\label{eq:faraday-rot}
\end{equation}
Status: \status{STANDARD\_EXACT} under the stated idealisation. Equation~\eqref{eq:faraday-rot} explains why a product containing \(B\omega N\) is natural in a rotating magnetic generator. It does \emph{not} make the source equation~\eqref{eq:source-zpe} identical to Faraday's law: Faraday contains a flux geometry and a derivative-generated trigonometric factor.

Maxwell's field-energy density supplies the ordinary electromagnetic baseline,
\begin{equation}
 u_{\rm EM}=\frac{\epsilon_0}{2}|\mathbf E|^2+\frac{1}{2\mu_0}|\mathbf B|^2.
\end{equation}
Any generator experiment must first close mechanical input, electromagnetic storage, ohmic loss, eddy currents, magnetic hysteresis, switching losses, and measured output before assigning a residual to a new channel.

\subsection{Quantum zero-point scale}

For a quantum harmonic oscillator,
\begin{equation}
E_n=\hbar\omega\left(n+\frac12\right),
\qquad
\boxed{E_0=\frac12\hbar\omega}.
\end{equation}
Status: \status{STANDARD\_EXACT}. This establishes a vacuum-ground-state energy scale in the oscillator model. It does not establish an extractable net-energy source. In this monograph, ``ZPE'' is therefore retained as source genealogy and as a research target, not as a closed energy-balance result.

\subsection{Dynamic phase and proper time}

For a local energy eigencomponent transported along proper time \(\tau\),
\begin{equation}
\boxed{d\varphi_{\rm dyn}=-\frac{E_{\rm loc}}{\hbar}\,d\tau.}
\label{eq:dynphase}
\end{equation}
For a worldline,
\begin{equation}
 d\tau=\frac1c\sqrt{-g_{\mu\nu}dx^\mu dx^\nu}
\end{equation}
for the mostly-plus convention. In a static metric one may write \(d\tau=N_g\,dt\), where \(N_g\) is the gravitational lapse for the chosen observer. The corresponding stationary redshift relation is conventionally represented as
\begin{equation}
 E_{\rm loc}=\frac{E_\infty}{N_g}.
\end{equation}
The invariant phase accounting prevents a redshift relation from becoming an energy-extraction rule by algebra alone.

PNCS v0.27 already contains a separate computational \emph{phase proper time}. Its current contract explicitly states that laboratory time is not modified. The present monograph preserves that boundary.

\subsection{Acceleration and the Unruh scale}

Uniform proper acceleration \(a\) defines the Unruh temperature
\begin{equation}
\boxed{T_U=\frac{\hbar a}{2\pi c k_B}},
\qquad
k_BT_U=\frac{\hbar a}{2\pi c}.
\end{equation}
Status: \status{STANDARD\_EXACT} in quantum field theory for the idealised uniformly accelerated detector setting \cite{Unruh1976}. The associated angular-frequency scale is
\begin{equation}
\omega_U=\frac{k_BT_U}{\hbar}=\frac{a}{2\pi c}.
\end{equation}
This provides a precise acceleration--frequency bridge. It is not a free-energy mechanism; the accelerated detector problem includes the work/interaction required to maintain acceleration.

\subsection{Aharonov--Bohm phase}

For a charged particle, the electromagnetic contribution to phase can be written, in a standard non-covariant split,
\begin{equation}
\boxed{
\Delta\varphi_{\rm AB}
=\frac{q}{\hbar}
\left(
\int \mathbf A\cdot d\mathbf l
-\int V\,dt
\right).
}
\end{equation}
For a closed magnetic Aharonov--Bohm loop,
\begin{equation}
\Delta\varphi_{\rm AB}=\frac{q}{\hbar}\oint\mathbf A\cdot d\mathbf l
=\frac{q\Phi_B}{\hbar}.
\end{equation}
The physical point is that a measurable interference phase can depend on gauge-potential holonomy even when the particle paths lie in a region of vanishing local magnetic field \cite{AharonovBohm1959}.

\subsection{Berry phase and geometric holonomy}

For a non-degenerate instantaneous eigenstate \(|n(\lambda)\rangle\), define the Berry connection
\begin{equation}
\mathcal A_a(\lambda)=i\langle n(\lambda)|\partial_{\lambda^a}n(\lambda)\rangle.
\end{equation}
For an adiabatically closed parameter loop \(C\),
\begin{equation}
\boxed{\gamma_B[C]=\oint_C\mathcal A_a\,d\lambda^a.}
\end{equation}
Its curvature is
\begin{equation}
\mathcal F=d\mathcal A,
\end{equation}
and, when the gauge patch permits,
\begin{equation}
\gamma_B=\int_\Sigma\mathcal F.
\end{equation}
Status: \status{STANDARD\_EXACT} for the usual adiabatic construction \cite{Berry1984}.

\subsection{Combined phase connection}

A convenient bookkeeping one-form is
\begin{equation}
\boxed{
\Gamma
= -\frac{E_{\rm loc}}{\hbar}d\tau
+\frac{q}{\hbar}A_\mu dx^\mu
+\mathcal A_a d\lambda^a,
}
\label{eq:combined-connection}
\end{equation}
with signs dependent on the electromagnetic four-potential convention. The corresponding holonomy is
\begin{equation}
\boxed{U[C]=\exp\left(i\oint_C\Gamma\right)}
\end{equation}
for an Abelian single-component case; non-Abelian transport requires path ordering.

\INTERP
Equation~\eqref{eq:combined-connection} motivates a PhaseNav principle: a state may carry physically or mathematically relevant information in the transport history, not only in a local endpoint coordinate. The interpretation is \emph{not} that gravitational, electromagnetic, Berry, and arithmetic connections are the same physical field. They are distinct connections that can share a transport/holonomy representation.

\subsection{Euler closure and the \((1,-1)\) sign pair}

Euler's identity gives
\begin{equation}
\boxed{e^{i\pi}=-1},\qquad e^{i2\pi}=+1.
\end{equation}
A loop with total phase \(\pi\pmod{2\pi}\) therefore returns the scalar holonomy \(-1\), while a double traversal gives \(+1\). Berry's original analysis includes sign reversal around certain degeneracies and relates geometric phase to familiar spinor sign behaviour \cite{Berry1984}.

\INTERP
The project uses the pair \((+1,-1)\) as a minimal two-valued holonomy representation. Whenever an arithmetic construction below exhibits one-loop \(-1\) and two-loop \(+1\), that is a representation-theoretic analogy to double-cover closure, not evidence that the arithmetic degree of freedom is physical spin.
